
\documentclass{lmcs}

\usepackage[T1]{fontenc}
\usepackage[utf8]{inputenc}
\usepackage{amssymb,mathtools,booktabs}
\microtypesetup{expansion=false}

\newcommand{\N}{\mathbb{N}}
\newcommand{\Z}{\mathbb{Z}}
\newcommand{\U}{\mathcal{U}}
\newcommand{\B}{\mathcal{B}}
\newcommand{\Lcal}{\mathcal{L}}
\newcommand{\Ocal}{\mathcal{O}}
\newcommand{\Gcal}{\mathcal{G}}
\newcommand{\Shape}{\Pi_{\!\infty}}
\newcommand{\Supp}{\mathrm{Supp}}
\newcommand{\hcomp}{\mathrm{hcomp}}
\newcommand{\compop}{\mathrm{comp}}
\newcommand{\transp}{\mathrm{transp}}
\newcommand{\Fib}{\mathrm{Fib}}

\title[Coherence Depth and Fibonacci Scaling]{Coherence Depth and Fibonacci Scaling in Intensional Type Theory}

\author[H.\ Lande]{Halvor Lande}
\address{Draft affiliation to be supplied}
\email{hsl@awc.no}

\keywords{intensional type theory, cubical type theory, coherence, univalence}
\subjclass[2012]{Theory of computation~$\rightarrow$~Type theory; Mathematics of computing~$\rightarrow$~Algebraic topology}

\begin{document}

\begin{abstract}
We study a metatheoretic invariant of framework extension in type theory. The paper models the worst-case scaling of \emph{foundational core extensions}: globally acting sealed additions such as universe hierarchies, global modalities, and heavily promoted generic interfaces whose public behavior must be specified against the active exported interface. We call the least historical depth at which the resulting obligations stabilize the coherence depth of the foundation. We formulate a payload-aware trace principle decomposing each exported interface into new core payload plus resolved coherence trace and, in the cubical case, a telescopic subsumption lemma showing that the depth-two arity bound is realized by a genuine two-layer chronological window. Together these turn coherence depth into an affine recurrence law for integration debt. In extensional systems with UIP the depth is one. In the fully coupled cubical univalent setting considered here the depth is two, and under uniform payload size the resulting integration law is shifted Fibonacci. We also give a Cubical Agda mechanization of both the depth theorem package itself and its affine recurrence consequence.
\end{abstract}

\maketitle

\section{Introduction}
\label{sec:introduction}

\subsection{The metatheoretic question}

This paper concerns a simple but surprisingly rigid question.

\begin{quote}
When a type-theoretic core calculus is extended by sealing a new structure against the existing interface, how does the cost of integration scale with the historical depth of the library?
\end{quote}

The question is not about transparent user developments inside a fixed library. It is about \emph{sealed} extensions: new type formers, modalities, or promoted interface packages whose public behavior becomes part of the global exported interface. More specifically, the paper studies the worst-case regime of \emph{foundational core extensions}: globally acting additions whose exported payload must be specified across the currently active interface. Examples include extending a univalent universe hierarchy, where the decoding machinery must be updated against the library's active codes; adding a global cohesive modality, which must distribute across the exported type formers and transports already present; and heavily promoted generic interfaces whose action is intended to range over the whole library. In that setting, extension is constrained not only by local typing rules, but also by the obligation to interact coherently with the eliminators, transports, computation principles, and comparison data already exported by the library.

Orthogonal extensions are not the subject of this recurrence law. Adjoining a new datatype such as \texttt{Tree} to a library that already contains \texttt{List} and \texttt{Socket} need not introduce primitive \texttt{Tree}--\texttt{Socket} interaction data merely to preserve canonicity. Such sparse or local extensions are important in practice, but they do not exhibit the maximal coupling pattern that this paper isolates. Ordinary user-level library growth therefore forms a sparse dependency DAG in general; the recurrence theorem studies only the asymptotic foundational-core regime of monolithic, globally acting extensions.

The central observation is that these obligations are stratified by historical depth. Some are purely unary: the new layer must act on one previously exported generator. Others are genuinely binary: the new layer must preserve a previously exported compatibility witness between two generators. The main invariant of the paper measures the least depth at which this process stabilizes.

\subsection{The core thesis}

We call that invariant the \emph{coherence depth} (or \emph{coherence window}) of the foundation. The thesis of the paper is that, once coherence depth is isolated, the asymptotic law of integration becomes forced.

\begin{quote}
In any fully coupled univalent core calculus with uniform payload size, framework extension debt exhibits Fibonacci scaling.
\end{quote}

Here ``fully coupled'' refers only to the foundational-core regime just described: admissible extensions are globally acting payloads whose semantics or promoted action is intended to range over the whole active interface, not orthogonal additions that leave most earlier exports untouched. Within that regime, every admissible extension must provide primitive interaction data against the whole active exported interface rather than against a chosen fragment, and maximal interface density is the canonicity-preserving form of that library-wide action. In extensional systems with UIP, the coherence depth collapses to~$1$, and under uniform payload size the recurrence is first-order affine. In the cubical univalent setting considered here, the coherence depth is~$2$, and after a constant payload shift the recurrence becomes Fibonacci.

\subsection{Summary of contributions}

The paper develops one theorem package with three components.

\begin{enumerate}
    \item We formalize \emph{core framework evolution}: states, sealed extensions, promoted interface packages, active interfaces, integration latency, and the distinction between orthogonal extensions and globally acting foundational core extensions. Within this framework we prove a maximal-interface-density theorem for the fully coupled foundational-core regime and a payload-aware \emph{Integration Trace Principle}, decomposing the exported interface of a sealed layer into its new public payload and the resolved obligations discharged in sealing it.
    \item We define coherence depth and prove a sharp dichotomy. In extensional type theory with UIP, the depth is~$1$. In the fully coupled cubical univalent setting, the depth is~$2$. The depth-two statement is supported by three convergent arguments: a dimensional upper bound from weak $\omega$-groupoid semantics and cubical Kan filling, a categorical lower bound from adjunction coherence, and a topological lower-bound family supplied by clutching constructions for principal bundles over spheres. The passage from the dimensional bound to a genuine two-step time window is made explicit by a telescopic subsumption lemma using cubical Kan filling.
    \item We combine the trace principle with the coherence-depth theorem to obtain a universal affine recurrence for integration latency. Under constant payload size~$c$, the depth-two case specializes to
    \[
      \Delta_{n+1} = \Delta_n + \Delta_{n-1} + 2c,
    \]
    which becomes homogeneous Fibonacci after the shift $U_n=\Delta_n+2c$. The associated structural inflation factor still converges to the Golden Ratio.
\end{enumerate}

The present version is a first full draft. The Cubical Agda development described in \autoref{sec:mechanization} now formalizes the opacity discipline, the payload-aware affine recurrence, the UIP collapse to depth~$1$, the algorithmic subsumption of arity-$3$ obligations by cubical composition, and the lower bound against global depth-one collapse. The dimensional, categorical, and topological arguments are still written here in condensed proof form because they explain the theorem architecture and its examples, but the paper's main depth claims no longer rest on prose alone.

\subsection{Related work}
\label{sec:related}

The higher-dimensional background comes from the interpretation of intensional identity types as weak $\omega$-groupoids, developed by Lumsdaine~\cite{lumsdaine} and van den Berg and Garner~\cite{vdberg-garner}. The operational setting for the intensional side of our depth dichotomy is cubical type theory in the sense of Cohen, Coquand, Huber, and M\"ortberg~\cite{cubical}, together with higher inductive types and Kan composition~\cite{cchm-hits}. For mechanization we use Cubical Agda~\cite{cubical-agda}.

The coherence-theoretic side is informed by Mac Lane's coherence theorem~\cite{maclane-coherence}. We also use Stasheff's analysis of higher associativity~\cite{stasheff} and the well-foundedness approach to coherence of Kraus and von Raumer~\cite{kraus-vonraumer}. Our contribution is to turn coherence from a qualitative structural property into a quantitative invariant of extension debt.

The topological family used for the lower bound is the classical clutching description of principal bundles over spheres; for standard bundle-theoretic background we refer to Husemoller~\cite{husemoller}. The univalent and cubical perspective sits in the general landscape of homotopy type theory and univalent foundations~\cite{hott}.

\section{Core framework evolution and canonicity}
\label{sec:framework}

We fix a foundation together with a discipline of \emph{sealed} extension. The library grows as a sequence
\[
  \B_0 \leadsto \B_1 \leadsto \B_2 \leadsto \cdots
\]
of states, where each step adjoins one new sealed layer.

\subsection{States, extensions, and promoted packages}

\begin{defi}[State]
\label{def:state}
A \emph{state} or \emph{library} is a monotone context $\B$ closed under derivability. An evolution step
\[
  \B_n \leadsto \B_{n+1}
\]
adjoins one sealed integration layer, denoted $L_{n+1}$.
\end{defi}

\begin{defi}[Candidate]
\label{def:candidate}
A \emph{candidate} over a library state $\B$ is a pair
\[
  X = (X_{\mathrm{core}}, \Gcal_{\mathrm{obl}}),
\]
where $X_{\mathrm{core}}$ is the mathematical content proposed for addition and $\Gcal_{\mathrm{obl}}$ is the family of atomic coherence obligations induced when $X$ is sealed against the active interface of~$\B$.
\end{defi}

\begin{defi}[Core payload size]
\label{def:payload}
For a candidate~$X$, let $P(X)$ be the finite set of public schemas contributed directly by $X_{\mathrm{core}}$ before any coherence witnesses against earlier layers are added. Its cardinality
\[
  \kappa(X) := |P(X)|
\]
is the \emph{core payload size} of~$X$.
\end{defi}

\begin{defi}[Sealed extension]
\label{def:sealed}
A \emph{sealed extension} is a candidate whose obligations have been discharged and whose resolved data are exported only through an opaque public interface. Future extensions may use those exports, but may not inspect the internal derivations that produced them.
\end{defi}

\begin{defi}[Promoted interface package]
\label{def:promoted}
A \emph{promoted interface package} over~$\B$ is a finite package of constructively irreducible generators and interaction data that is exported as one new sealed layer. Such a package need not modify the reduction rules of the kernel. What matters is that its public generators become part of the active exported interface for later sealing steps.
\end{defi}

\begin{defi}[Integration latency]
\label{def:latency}
The \emph{integration latency} of a candidate $X$ over~$\B$ is
\[
  \Delta(X \mid \B) := |\Gcal_{\mathrm{obl}}|,
\]
the number of atomic obligations that must be resolved before $X$ can be exported as a sealed layer.
\end{defi}

For a sealed layer $L_k$ obtained from a candidate~$X_k$, we write $S(L_k)$ for its set of exported interface schemas and
\[
  \kappa_k := \kappa(X_k)
\]
for its core payload size. The indexing is chosen so that $\Delta_k$ denotes the sealing cost of~$L_k$.

\begin{prop}[Transparent user-level code lies outside the recurrence]
\label{prop:transparent}
Let $U$ be a family of transparent definitions, lemmas, or record packages over~$\B$ whose symbols are definitionally reducible using the existing operational semantics and are not promoted to a new opaque exported layer. Then $U$ contributes zero integration latency and does not alter the active interface seen by later sealing steps.
\end{prop}

\begin{proof}
Such a development creates no new irreducible interaction site against the exported interface. All reductions remain inherited from the pre-existing library, and no new opaque generators are added to the global interface telescope. Hence $U$ is ordinary growth inside one state, not a state transition of the form $\B_n \leadsto \B_{n+1}$.
\end{proof}

\begin{defi}[Foundational core extension]
\label{def:foundational-core}
A candidate~$X$ over~$\B_n$ is a \emph{foundational core extension} if its payload is intended to act globally on the currently active interface: each public generator in $P(X)$ is operationally or semantically incomplete until its behavior has been specified against every generator in the designated coupling footprint inside the active exported interface. Typical examples are extensions of a universe-decoding apparatus, global modalities whose action must distribute across existing type formers and transports, and heavily promoted generic interfaces whose exported action is library-wide.
\end{defi}

\begin{rem}
\label{rem:orthogonal}
Not every sealed extension is foundational in the sense of \autoref{def:foundational-core}. A new datatype or API layer may be orthogonal to much of the existing library and therefore require only sparse interaction data. Such extensions may preserve canonicity without becoming globally coupled to earlier layers. The recurrence theorem of this paper is intentionally a worst-case statement about foundational core extensions, not a universality claim about all sealed library growth.
\end{rem}

This scope restriction is essential for the density theorem below. Ordinary user-level additions typically enlarge a sparse DAG of modules or datatypes and do not force primitive interaction with unrelated earlier exports. The maximal-coupling statement proved here isolates the special foundational-core regime in which the new sealed layer is advertised as a \emph{global} operation on the active interface, so that canonicity and global closure become exhaustiveness requirements rather than bookkeeping conventions.

\subsection{Maximal interface density as a canonicity requirement}

To speak about coupling to historical layers, it is convenient to parameterize the visible interface by depth.

\begin{defi}[Historical interface]
\label{def:historical-interface}
For $k \geq 1$, the \emph{historical interface of depth~$k$} at stage~$n$ is
\[
  I_n^{(k)} := \biguplus_{j=0}^{k-1} S(L_{n-j}),
\]
the tagged coproduct of the schemas exported by the most recent $k$ sealed layers.
\end{defi}

\begin{defi}[Maximal interface density]
\label{def:max-density}
A foundational core candidate sealed against $I_n^{(k)}$ has \emph{maximal interface density} if every generator exported by $I_n^{(k)}$ lies in its coupling footprint and receives exactly one primitive interaction datum from the new extension. For native type formers such a datum is a defining reduction clause or decoding equation; for global modalities it is a distribution or transport law; for promoted interface packages it is a generatorwise comparison, transport, or action witness.
\end{defi}

\begin{defi}[Fully coupled foundation]
\label{def:fully-coupled}
A foundation is \emph{fully coupled} if the admissible sealed extensions under study are foundational core extensions and each such extension is required to satisfy maximal interface density with respect to the active historical interface. Equivalently, the paper's admissible extensions in a fully coupled foundation are globally acting payloads whose coupling footprint is the whole active interface, not orthogonal additions with sparse dependency patterns.
\end{defi}

The next theorem turns maximal interface density into a metatheoretic necessity for this globally coupled regime rather than a stylistic choice.

\begin{thm}[Canonicity forces maximal interface density for foundational core extensions]
\label{thm:canonicity}
Let $\B_n$ be a library state in a fully coupled foundation, and let $X$ be an admissible foundational core extension sealed against the historical interface $I_n^{(k)}$. If the extended theory is to preserve the global closure and canonicity discipline of the sealed library, then every generator of $I_n^{(k)}$ must receive exactly one primitive interaction datum from~$X$. Consequently,
\[
  \Delta(X \mid \B_n) = |I_n^{(k)}|.
\]
\end{thm}

\begin{proof}
By \autoref{def:fully-coupled}, the admissible candidates considered here are not arbitrary sealed additions but foundational core extensions whose payload is intended to act on the whole active interface. The theorem is therefore about that global regime only; orthogonal extensions of the kind noted in \autoref{rem:orthogonal} are outside its hypotheses.

First consider a native global computational extension. In a Tarski-style universe presentation, suppose the active interface exports a code universe $U$ and a decoding map $\mathsf{El} : U \to \mathsf{Type}$. If the new sealed layer extends the universe with a fresh code but the extended decoding machinery is not given the corresponding new primitive computation clause, then closed terms mentioning that code no longer reduce through $\mathsf{El}$, so canonicity fails on the extended language. More generally, if a new global operator such as a modality, transport schema, or promoted eliminator package is advertised as acting on the whole active interface but omits some generator of $I_n^{(k)}$, then the extended reduction or comparison relation is partial on well-typed terms built from that generator. In compiler terms, the global action has failed the relevant exhaustiveness condition.

For globally promoted interface packages the same point appears at the level of exported semantics. If the new opaque package is meant to provide a library-wide action but omits the comparison witness for some previously exported generator, then later sealing steps cannot treat it as a closed global operation ranging over the whole library. The public action is only partially defined on the active interface, so the sealed calculus is not globally closed.

This proves the ``at least one'' direction for foundational core extensions. For the converse, the issue is not metaphysical uniqueness of higher structure but the single-clause discipline required for computation. For native global extensions, two distinct primitive defining clauses for the same head and the same prior generator would overlap and destroy confluence. For globally promoted packages, once the public source, target, and prior generator are fixed, a second primitive datum would advertise competing basic behavior for the same global action. Hence exactly one primitive clause is forced at each head/generator site, and the number of obligations is precisely the cardinality of the historical interface.
\end{proof}

\subsection{The integration trace principle}

In the fully coupled foundational-core regime, the preceding theorem identifies sealing cost with the size of the active interface. The next theorem explains how those resolved obligations reappear in the next layer's public export alongside the layer's own new mathematical payload.

\begin{thm}[Integration Trace Principle]
\label{thm:trace}
Each sealed integration layer exports a disjoint union of two kinds of public schema: the new payload contributed by its core content and the trace schemas corresponding to its resolved obligations. Equivalently, if $L_k$ is the layer obtained by sealing a candidate of core payload size~$\kappa_k$ and coherence cost~$\Delta_k$, then there is a canonical decomposition
\[
  S(L_k) \cong S_{\mathrm{core}}(L_k) \uplus S_{\mathrm{tr}}(L_k)
\]
with
\[
  |S_{\mathrm{core}}(L_k)| = \kappa_k,
  \qquad
  |S_{\mathrm{tr}}(L_k)| = \Delta_k.
\]
Hence
\[
  |S(L_k)| = \kappa_k + \Delta_k.
\]
\end{thm}

\begin{proof}
There are two logically distinct sources of exported schemas.

First, the candidate's core content $X_{k,\mathrm{core}}$ contributes genuinely new public generators: a new type former, modality, promoted interface operator, or other primitive schema that did not previously exist in the library. By definition of core payload size, the number of these newly exported payload schemas is~$\kappa_k$. These form $S_{\mathrm{core}}(L_k)$.

Second, by \autoref{thm:canonicity}, sealing against the active interface requires exactly one primitive interaction datum for each active exported generator. Once such a datum is supplied and the layer is sealed, the datum no longer survives as an internal derivation visible to future extensions. It survives only through the public schema recording how the new payload acts on that older generator. These schemas form $S_{\mathrm{tr}}(L_k)$, and there is one of them for each resolved obligation. Hence $|S_{\mathrm{tr}}(L_k)|=\Delta_k$.

Because the layer is sealed, every public schema of $L_k$ is of one of these two kinds: either it is a newly introduced payload symbol, or it is the exported trace of a resolved obligation attaching that payload to prior interface data. The two kinds are disjoint, since payload symbols are primitive generators introduced by the new layer itself, whereas trace schemas are indexed by the prior interface against which the new layer was sealed. Therefore the public export decomposes canonically as
\[
  S(L_k) \cong S_{\mathrm{core}}(L_k) \uplus S_{\mathrm{tr}}(L_k),
\]
and the cardinality formula follows immediately.
\end{proof}

\begin{rem}
\label{rem:trace-native-vs-api}
The trace principle counts kernel-level extensions and promoted interface packages uniformly. In both cases the public export has a payload part and a trace part. The only difference is the kind of primitive datum involved: for native extensions the payload contributes new reduction heads and the trace contributes defining clauses, whereas for promoted packages the payload contributes new exported operators and the trace contributes comparison or transport witnesses.
\end{rem}

\section{The coherence window concept and the depth dichotomy}
\label{sec:coherence}

\subsection{Induced obligations and historical arity}

\begin{defi}[Induced obligations]
\label{def:obligations}
Let $\Ocal^{(k)}(X)$ be the set of normalized atomic obligations induced when candidate~$X$ is sealed against a historical interface of depth~$k$. Since the interface grows cumulatively, there are inclusions
\[
  \Ocal^{(1)}(X) \subseteq \Ocal^{(2)}(X) \subseteq \Ocal^{(3)}(X) \subseteq \cdots.
\]
\end{defi}

\begin{defi}[Historical support and arity]
\label{def:arity}
For a normalized obligation~$o$, let $\Supp(o)$ be the set of sealed layers whose exported schemas occur irreducibly in the typing derivation of~$o$ after each prior layer has been collapsed to its opaque interface. Its \emph{historical arity} is
\[
  a(o) := |\Supp(o)|.
\]
\end{defi}

\begin{rem}
\label{rem:deep-lookups}
Historical arity counts independent prior layers, not textual age. A direct lookup to an old exported theorem is still unary, because it lands in one opaque interface. Multi-layer debt arises only when the new layer must cohere simultaneously with several independent prior exports.
\end{rem}

\begin{lem}[Arity-to-dimension correspondence]
\label{lem:arity-dimension}
An irreducible obligation of historical arity~$k$ requires a $k$-dimensional coherence cell. Arity~$1$ produces substitution or transport data, arity~$2$ produces a naturality square or homotopy, and arity $k \geq 3$ would require a genuine higher filler.
\end{lem}

\begin{proof}
With one historical input, the new layer must supply a comparison morphism between its behavior and one prior export, so the obligation is one-dimensional. With two independent historical inputs, the witness compares two unary composites and therefore has the form of a square or homotopy. More generally, relating $k$ independently supplied historical interfaces requires a $k$-dimensional filler between the corresponding composites.
\end{proof}

\subsection{Definition of coherence depth}

\begin{defi}[Coherence depth]
\label{def:depth}
A foundation has \emph{coherence depth}~$d$ if $d$ is the least integer such that for every candidate~$X$ and all $k \geq d$ there is a canonical equivalence
\[
  \Ocal^{(k)}(X) \simeq \Ocal^{(d)}(X).
\]
We also call~$d$ the \emph{coherence window} of the foundation.
\end{defi}

The point of \autoref{def:depth} is that beyond depth~$d$, deeper historical interaction contributes no new independent obligations. Once that stabilization has been upgraded to a genuine chronological window, the size of the primitive active interface controls the recurrence law for integration debt.

\begin{rem}
\label{rem:arity-vs-chronology}
\autoref{def:depth} is a stabilization statement about \emph{obligation sets}. It should not be read as saying that a bound on coherence dimension automatically identifies the chronological layers that must still be inspected primitively. In particular, an arity-two obligation could in principle couple a new layer to a remote historical layer. To extract a genuine time window from a dimensional bound one needs an additional subsumption argument showing that such remote binary comparisons are already generated by recently exported traces. In the cubical univalent case, this passage from arity to chronology is supplied by \autoref{lem:telescopic}.
\end{rem}

\subsection{\texorpdfstring{Theorem A: extensional systems collapse to $d=1$}{Theorem A: extensional systems collapse to d=1}}

\begin{thm}
\label{thm:extensional}
In Martin-L\"of type theory with UIP, the coherence depth is~$1$.
\end{thm}

\begin{proof}
Under UIP, any two proofs of identity between fixed endpoints are equal. Hence every identity type is a set, and every higher identity type is a proposition; once inhabited, it is contractible. Consider any normalized coherence obligation of dimension at least~$2$. By \autoref{lem:arity-dimension}, it corresponds to a higher identity witness. But in the presence of UIP such a witness is uniquely determined by its boundary and therefore carries no independent data. Thus every obligation of arity at least~$2$ collapses to data already fixed by its unary boundary. Equivalently,
\[
  \Ocal^{(k)}(X) \simeq \Ocal^{(1)}(X)
  \qquad
  (k \geq 1).
\]
Since depth~$0$ would mean the absence of any interaction with the current interface, which is impossible for a nontrivial sealing step, the least stabilizing depth is~$1$.
\end{proof}

\subsection{\texorpdfstring{Theorem B: fully coupled univalent systems have $d=2$}{Theorem B: fully coupled univalent systems have d=2}}

From now on we work in a fully coupled cubical setting in the style of Cohen, Coquand, Huber, and M\"ortberg~\cite{cubical,cchm-hits}, with univalence and Kan composition available in the ambient operational semantics. The next three ingredients are deliberately complementary: the dimensional argument gives the upper bound, while the categorical and topological families give independent lower bounds. Together they pin down exact depth~$2$, and the telescopic lemma makes explicit how the dimensional bound is realized as a genuine two-step time window.

One semantic clarification is important before the cubical argument begins. In univalent foundations, types behave as weak $\omega$-groupoids, so the geometric space of fillers for a fixed higher boundary is not generally contractible; open $3$-boxes can admit many fillers. The recurrence analyzed in this paper does not count that homotopy-theoretic space. It counts only \emph{primitive integration data}: the new opaque generators and defining clauses that a developer must add at sealing time. In cubical type theory, the operational semantics provide generic Kan operations such as $\hcomp$, $\compop$, and $\transp$ that already inhabit many higher boundaries uniformly. The crucial distinction is therefore syntactic: a higher witness contributes to the recurrence only when it must be entered as new domain-specific primitive payload, not when it is synthesized by the ambient calculus.

\begin{defi}[Primitive cost of an atomic obligation]
\label{def:primitive-cost}
Let $o$ be a normalized atomic obligation. We set $\delta(o)=1$ when resolving $o$ requires the candidate being sealed to contribute a new domain-specific opaque generator, defining equation, or pattern-matching clause. We set $\delta(o)=0$ when a witness for $o$ is uniformly synthesized from already exported boundary data using only generic framework term formers.

An obligation is \emph{irreducible} exactly when $\delta(o)=1$. Thus the recurrence for $\Delta_n$ counts primitive obligations only; a zero-cost witness is derived trace rather than new sealing payload.
\end{defi}

\paragraph{A concrete cubical sketch.}
The abstract distinction between payload and trace can be made visible in a tiny Agda-shaped example without importing the full PEN framework. Suppose an older sealed layer exports a small Tarski universe with function codes and their decoding equation. A new cohesive layer contributes a genuinely new modal payload \texttt{Flat} together with its unit \texttt{eta}. Everything else needed to make that modality act on the older universe is integration trace: a lifted action on old codes, a comparison between decoding before or after lifting, and the distribution law over old function codes. In the snippet below, the witness \texttt{remote-face} is the point of the example. It is not additional payload. It is exactly the missing comparison face of the open $3$-box determined by the old decoding law together with the new unary and binary trace data, and cubical composition computes it uniformly.

\begin{quote}
\begin{verbatim}
-- Schematic Cubical Agda: old universe + new payload + induced trace

record OldU : Set1 where
  field
    U      : Set
    El     : U -> Set
    Arr    : U -> U -> U
    El-Arr : (A B : U) ->
      Path Set (El (Arr A B)) (El A -> El B)

record Payload : Set1 where
  field
    Flat : Set -> Set
    eta  : {A : Set} -> A -> Flat A

record Trace (L : OldU) (P : Payload) : Set1 where
  open OldU L
  open Payload P
  field
    FlatU         : U -> U
    flat-El       : (A : U) ->
      Path Set (El (FlatU A)) (Flat (El A))
    flat-code-Arr : (A B : U) ->
      Path U (FlatU (Arr A B)) (Arr (FlatU A) (FlatU B))
    flat-Arr      : (A B : U) ->
      Path Set (Flat (El A -> El B))
               (Flat (El A) -> Flat (El B))

module _ (L : OldU) (P : Payload) (T : Trace L P) where
  open OldU L
  open Payload P
  open Trace T

  -- left route:
  --   El (FlatU (Arr A B))
  --     -> Flat (El (Arr A B))
  --     -> Flat (El A -> El B)
  --     -> Flat (El A) -> Flat (El B)

  -- right route:
  --   El (FlatU (Arr A B))
  --     -> El (Arr (FlatU A) (FlatU B))
  --     -> El (FlatU A) -> El (FlatU B)
  --     -> Flat (El A) -> Flat (El B)

  postulate
    left-route  :
      (A B : U) ->
      Path Set (El (FlatU (Arr A B))) (Flat (El A) -> Flat (El B))
    right-route :
      (A B : U) ->
      Path Set (El (FlatU (Arr A B))) (Flat (El A) -> Flat (El B))

  -- The open 3-box says these two boundary composites agree.
  -- This face is derived trace, not new payload.
  remote-face :
    (A B : U) ->
    Path (Path Set (El (FlatU (Arr A B)))
                   (Flat (El A) -> Flat (El B)))
         (left-route A B) (right-route A B)
  remote-face A B = hcomp ...
\end{verbatim}
\end{quote}

Read top-to-bottom, \texttt{Payload} is the human-authored new content, while \texttt{Trace} is the compatibility debt forced by sealing against the old universe. The two routes are the two composites from \texttt{El (FlatU (Arr A B))} to \texttt{Flat (El A) -> Flat (El B)}. The arity-$3$ obligation is precisely the comparison between those composites. In a cubical presentation that comparison is not another primitive clause that must be postulated separately for each old constructor; it is the filler of the open box generated by the already exported boundary.

\subsubsection{\texorpdfstring{Dimensional upper bound ($d \leq 2$)}{Dimensional upper bound}}

\begin{thm}[Dimensional upper bound]
\label{thm:upper}
In the fully coupled cubical univalent setting, every irreducible coherence obligation has historical arity at most~$2$.
\end{thm}

\begin{proof}
Let $o$ be a normalized atomic obligation of historical arity $k \geq 3$. By \autoref{lem:arity-dimension}, its witness has the form of a $k$-dimensional coherence filler: syntactically, one must inhabit a higher identity type whose boundary is an open $k$-box assembled from lower-dimensional traces already exported by prior layers.

In a generic intensional theory, such a filler could require an additional postulate or ad hoc clause. In that situation one would have $\delta(o)=1$ in the sense of \autoref{def:primitive-cost}. In the cubical setting considered here, however, Kan composition is part of the ambient syntax. For every well-typed open box boundary of dimension at least~$3$, the term former $\hcomp\, r\, u\, u_0$ computes a canonical inhabitant from the boundary system~$u$ and cap~$u_0$ using only the generic typing rules of the calculus~\cite{cubical,cchm-hits}. No new domain-specific opaque generator, defining equation, or pattern-matching clause is introduced.

Hence every arity-$k$ obligation with $k \geq 3$ has primitive cost $\delta(o)=0$. Such an obligation may still have a nontrivial homotopy-theoretic space of fillers, but it contributes no primitive sealing payload and is therefore not irreducible. By contrast, unary and binary obligations may still require genuine developer-supplied data. Therefore every irreducible coherence obligation has historical arity at most~$2$.
\end{proof}

\begin{lem}[Telescopic subsumption]
\label{lem:telescopic}
In the fully coupled cubical univalent setting, primitive sealing against the two most recent layers algorithmically synthesizes all older binary comparisons. More precisely, let $X$ be a candidate added at stage~$n+1$. Once $X$ has supplied its primitive unary and binary data against $L_n$ and $L_{n-1}$, every would-be primitive comparison between $X$ and a layer $L_{n-r}$ with $r \geq 2$ is uniformly derived by iterated cubical Kan filling from the traces already exported by $L_n,\dots,L_{n-r+1}$.
\end{lem}

\begin{proof}
Start with the first nonadjacent layer $L_{n-2}$. Because $L_n$ was itself sealed against $L_{n-1}$ and $L_{n-2}$, its exported trace includes a unary comparison
\[
  f_n : L_n \to L_{n-1},
\]
the previous layer exports
\[
  g_{n-1} : L_{n-1} \to L_{n-2},
\]
and the binary part of the trace of $L_n$ contains a $2$-cell
\[
  \alpha_n : g_{n-1}\circ f_n \Rightarrow h_n
\]
recording the induced comparison $h_n : L_n \to L_{n-2}$. When $X$ is sealed against the last two layers, it supplies unary comparisons
\[
  x_n : X \to L_n,
  \qquad
  x_{n-1} : X \to L_{n-1},
\]
and a binary coherence square
\[
  \beta_n : f_n \circ x_n \Rightarrow x_{n-1}.
\]
The cells $x_n$, $x_{n-1}$, $f_n$, $g_{n-1}$, $\alpha_n$, and $\beta_n$, together with the evident degenerate composition faces, determine an open $3$-box. Its missing face is precisely a comparison
\[
  \gamma_{n-2} : h_n \circ x_n \Rightarrow g_{n-1}\circ x_{n-1}
\]
between the two composites from $X$ to $L_{n-2}$. By \autoref{def:primitive-cost}, this comparison has primitive cost~$0$: cubical Kan filling via $\hcomp$ computes a canonical filler from the displayed boundary data, so no additional domain-specific clause is added when sealing~$X$.

The same argument now iterates. Suppose inductively that a comparison from $X$ to $L_{n-r+1}$ has already been syntactically derived. Since $L_{n-r+1}$ was sealed against $L_{n-r}$ and $L_{n-r-1}$, its exported trace again contains the adjacent unary bridge to $L_{n-r}$ together with the corresponding binary square witnessing the comparison of composites. Combining that previously exported square with the already derived comparison for $X$ forms another open $3$-box, and $\hcomp$ fills it again. The resulting derived face supplies the comparison from $X$ to $L_{n-r}$. At each additional historical step, the older layer contributes boundary data only; it contributes no new primitive cell. Thus layers at distance at least two are telescopically subsumed by the adjacent traces already exported by the recent history.
\end{proof}

\begin{cor}[Two-layer chronological window]
\label{cor:chrono-window}
In the fully coupled cubical univalent setting, every sealing step at stage~$n+1$ has a genuine chronological window of size~$2$: all primitive obligations are generated by the exported traces of $L_n$ and $L_{n-1}$, while comparisons with layers $L_{n-r}$ for $r \geq 2$ are algorithmically synthesized by iterated Kan filling. Equivalently, the pair $(L_n,L_{n-1})$ forms a chronological Markov blanket for subsequent sealing.
\end{cor}

\begin{proof}
\autoref{thm:upper} shows that no irreducible obligation has historical arity greater than~$2$. \autoref{lem:telescopic} upgrades that dimensional statement to a temporal one: any binary comparison with a remote layer is induced from the adjacent trace data already exported by the last two layers. Hence primitive sealing data are exhausted by $S(L_n)\uplus S(L_{n-1})$, and all deeper history contributes only derived fillers.
\end{proof}

\subsubsection{\texorpdfstring{Categorical lower bound: the adjunction barrier ($d \geq 2$)}{Categorical lower bound}}

\begin{thm}[Adjunction barrier]
\label{thm:adjunction}
A depth-one system cannot verify the genuine coherence of a nontrivial adjunction. Consequently, any fully coupled univalent setting supporting adjoint structure has coherence depth at least~$2$.
\end{thm}

\begin{proof}
Assume for contradiction that the ambient univalent foundation collapses to depth~$1$. By \autoref{def:depth}, this means that every binary obligation is already determined by unary boundary data:
\[
  \Ocal^{(2)}(X) \simeq \Ocal^{(1)}(X)
  \qquad
  \text{for all candidates } X.
\]
Applied to identity-type coherence, this yields the global contraction principle formalized in \path{Metatheory/AdjunctionBarrier.agda}:
\[
  (A\, B : \U)\,(p\, q : A \equiv B) \to \mathsf{isContr}(p \equiv q).
\]
Equivalently, once the endpoints are fixed, any two witnesses of identity are propositionally equal; this is the relevant UIP-style collapse of binary coherence spaces.

Now take the two-point type $\mathbf{2}$ and the nontrivial auto-equivalence
\[
  \mathsf{swap} : \mathbf{2} \simeq \mathbf{2}
\]
that exchanges the two points. Univalence provides a path
\[
  p := \mathsf{ua}(\mathsf{swap}) : \mathbf{2} \equiv \mathbf{2}.
\]
If depth~$1$ held, then the binary obligation space between $\mathsf{refl}_{\mathbf{2}}$ and~$p$ would be contractible, so in particular we would obtain
\[
  p \equiv \mathsf{refl}_{\mathbf{2}}.
\]
But transport along $\mathsf{ua}(\mathsf{swap})$ computes as the action of $\mathsf{swap}$, whereas transport along $\mathsf{refl}$ computes as the identity. Applying the displayed equality to the point $\mathsf{left} : \mathbf{2}$ therefore forces $\mathsf{swap}(\mathsf{left})=\mathsf{left}$, contradicting $\mathsf{swap}(\mathsf{left})=\mathsf{right}$ and $\mathsf{left}\neq\mathsf{right}$. This is exactly the contradiction formalized in the Agda development by proving that \texttt{swap-path} is nontrivial and then deriving \texttt{depth1-insufficient}.

Therefore a univalent system cannot collapse globally to depth~$1$. Binary coherence spaces do not all contract, so a fully coupled setting cannot derive every $2$-cell from unary boundary data alone. In particular, the triangle identities of an adjunction remain genuine binary structure. An adjunction $L \dashv R$ is specified by functors $L$ and $R$, together with a unit and counit
\[
  \eta : 1 \to RL,
  \qquad
  \varepsilon : LR \to 1,
\]
satisfying the triangle identities
\[
  \varepsilon L \circ L\eta \simeq \mathrm{id}_L,
  \qquad
  R\varepsilon \circ \eta R \simeq \mathrm{id}_R.
\]
The triangle identities are genuine $2$-cells, and the preceding contradiction shows that such binary coherence cannot be uniformly collapsed away in the univalent setting. Hence depth one is insufficient, and $d \geq 2$.
\end{proof}

\subsubsection{Topological exact family: clutching constructions}

\begin{thm}[Clutching family]
\label{thm:clutching}
The HIT clutching presentation over $S^2$ furnishes an exact depth-$2$ witness: a genuine binary comparison is necessary to specify the bundle data, but no independent primitive arity-$3$ obligation occurs.
\end{thm}

\begin{proof}
We translate the clutching construction into higher-inductive syntax. In homotopy type theory, the $2$-sphere is the suspension $S^2 := \Sigma S^1$, generated by point constructors
\[
  \mathsf{N}, \mathsf{S} : S^2
\]
and a path constructor
\[
  \mathsf{merid} : S^1 \to (\mathsf{N} \equiv \mathsf{S})
\]
~\cite{hott,cchm-hits}. To specify an exported layer whose total space is a principal bundle over~$S^2$, it suffices to define a type family
\[
  P : S^2 \to \U.
\]
By the dependent eliminator for the suspension, such a family is determined by exactly three pieces of data:
\[
  P_{\mathsf{N}} : \U,
  \qquad
  P_{\mathsf{S}} : \U,
  \qquad
  \gamma : \Pi_{x:S^1} \bigl(P_{\mathsf{N}} \equiv P_{\mathsf{S}}\bigr).
\]
By univalence, the last component may equally be read as a family of equivalences between the two pole fibers, which is the type-theoretic clutching datum.

Now $S^1$ is itself a higher inductive type generated by a point $\mathsf{base}$ and a loop $\mathsf{loop} : \mathsf{base} \equiv \mathsf{base}$. Therefore giving $\gamma$ requires exactly two layers of information: first a base equivalence
\[
  \gamma(\mathsf{base}) : P_{\mathsf{N}} \equiv P_{\mathsf{S}},
\]
which is unary data, and second a path over $\mathsf{loop}$ identifying the two transports induced by that equivalence around the circle. That second datum is a genuine binary coherence cell: it compares two unary composites and cannot be recovered from the endpoints alone. Hence depth~$1$ does not suffice.

Crucially, the HIT signature of $S^2$ contains no constructor of dimension~$3$ or higher. Accordingly, the elimination rule for $S^2$ asks for no primitive input beyond the two point data and the path family~$\gamma$ above. Once the binary path-over-loop has been supplied, any further compatibility condition is generated by the ordinary computation and transport rules of the eliminator, not by a new developer-supplied clause. In the sense of \autoref{def:primitive-cost}, every would-be arity-$3$ condition has cost~$0$. Thus the clutching presentation furnishes an exact depth-$2$ witness.
\end{proof}

\begin{cor}
\label{cor:d2}
In the fully coupled cubical univalent setting considered above, the coherence depth is exactly~$2$.
\end{cor}

\begin{proof}
\autoref{thm:upper} provides the upper bound by showing that no independent obligation survives beyond the $2$-skeleton. \autoref{lem:telescopic} shows that this depth-two upper bound is realized by a genuine two-layer chronological window rather than merely by a bound on cell dimension. \autoref{thm:adjunction} and \autoref{thm:clutching} provide lower bounds by exhibiting genuine depth-two phenomena that do not collapse in the univalent setting. Therefore the least stabilizing depth is~$2$.
\end{proof}

\section{The complexity scaling theorem}
\label{sec:scaling}

Let $L_n$ denote the $n$th sealed layer and write
\[
  \Delta_n := \Delta(L_n \mid \B_{n-1})
\]
for its sealing cost, and let $\kappa_n$ be its core payload size. The recurrence theorem counts \emph{primitive} sealing data, so it uses not only the stabilization statement of \autoref{def:depth} but also a theorem identifying a genuine chronological window. In the depth-two cubical univalent case, this is exactly \autoref{cor:chrono-window}; in the UIP case, the proof of \autoref{thm:extensional} shows that all non-unary comparisons collapse to their boundary, so the window is one layer. More generally, whenever a foundation has coherence depth~$d$ realized by a chronological window of size~$d$, the primitive active interface at stage~$n+1$ is the depth-$d$ historical interface built from the \emph{full} exports of the preceding layers:
\[
  I_n^{(d)} = \biguplus_{j=0}^{d-1} S(L_{n-j}).
\]
For the initial bootstrap stages before the full window has been populated, we adopt the natural convention that nonexistent earlier layers contribute the empty interface.

\subsection{The universal affine recurrence}

\begin{thm}[Universal affine recurrence]
\label{thm:recurrence}
For every foundation whose coherence depth~$d$ is realized by a chronological window of size~$d$, the integration latency satisfies
\[
  \Delta_{n+1} = \sum_{j=0}^{d-1} \bigl(\Delta_{n-j} + \kappa_{n-j}\bigr)
  \qquad
  (n \geq d).
\]
\end{thm}

\begin{proof}
By definition, the next sealing cost is the size of the active interface:
\[
  \Delta_{n+1} = |I_n^{(d)}|.
\]
Using the tagged coproduct decomposition of \autoref{def:historical-interface},
\[
  |I_n^{(d)}|
  =
  \sum_{j=0}^{d-1} |S(L_{n-j})|.
\]
Now apply the Integration Trace Principle, \autoref{thm:trace}, which gives
\[
  |S(L_{n-j})| = \Delta_{n-j} + \kappa_{n-j}
\]
for each~$j$. Substituting yields the claimed affine recurrence.
\end{proof}

For later comparison, write
\[
  \tau_n := \sum_{i=1}^n \Delta_i
\]
for the cumulative integration latency through stage~$n$.

\subsection{\texorpdfstring{Affine one-step growth at $d=1$}{Affine one-step growth at d=1}}

\begin{cor}
\label{cor:d1}
If the coherence depth is~$1$, this depth is realized by a one-layer chronological window, and every layer contributes the same core payload size~$c$, then
\[
  \Delta_{n+1} = \Delta_n + c.
\]
If moreover the initial state exports no prior sealed layers, then
\[
  \Delta_1 = 0,
  \qquad
  \Delta_n = (n-1)c,
\]
and therefore
\[
  \tau_n = c\,\frac{n(n-1)}{2}.
\]
In particular, integration latency grows linearly and cumulative integration cost grows quadratically.
\end{cor}

\begin{proof}
Substitute $d=1$ into \autoref{thm:recurrence}. Then only the term $j=0$ remains, so
\[
  \Delta_{n+1} = \Delta_n + \kappa_n.
\]
Under the uniform payload assumption $\kappa_n=c$, this becomes $\Delta_{n+1}=\Delta_n+c$. If the initial state exports no prior sealed layers, then the first sealing step sees the empty interface and therefore has cost $\Delta_1=0$. Induction gives $\Delta_n=(n-1)c$, and summing this arithmetic progression yields the stated formula for~$\tau_n$.
\end{proof}

\subsection{\texorpdfstring{Affine Fibonacci scaling at $d=2$}{Affine Fibonacci scaling at d=2}}

By \autoref{cor:chrono-window}, the fully coupled cubical univalent depth-two theorem is chronologically realized by the two most recent sealed traces.

Assume in this subsection that every sealed layer contributes the same core payload size~$c$, and write $(F_n)_{n \geq 1}$ for the Fibonacci sequence with $F_1=F_2=1$.

\begin{cor}[Affine Fibonacci recurrence]
\label{cor:fibonacci}
In the fully coupled cubical univalent setting, if the coherence depth is~$2$ and every sealed layer contributes core payload size~$c$, then
\[
  \Delta_{n+1} = \Delta_n + \Delta_{n-1} + 2c.
\]
If we define
\[
  U_n := \Delta_n + 2c,
\]
then
\[
  U_{n+1} = U_n + U_{n-1}.
\]
If moreover the initial state exports no prior sealed layers, then
\[
  \Delta_1 = 0,
  \qquad
  \Delta_2 = c,
\]
so
\[
  U_n = cF_{n+2},
  \qquad
  \Delta_n = c(F_{n+2}-2),
\]
and the cumulative latency satisfies
\[
  \tau_n = c(F_{n+4}-2n-3).
\]
\end{cor}

\begin{proof}
The recurrence is the case $d=2$ of \autoref{thm:recurrence} together with the uniform payload assumption $\kappa_n=c$, giving
\[
  \Delta_{n+1} = \Delta_n + \Delta_{n-1} + 2c.
\]
Adding $2c$ to both sides yields
\[
  U_{n+1} = U_n + U_{n-1}.
\]

If the initial state exports no prior sealed layers, then the first layer has no earlier interface against which to cohere, so $\Delta_1=0$. The second layer sees only the first layer's export, whose size is $\kappa_1+\Delta_1=c$, hence $\Delta_2=c$. Therefore
\[
  U_1 = 2c = cF_3,
  \qquad
  U_2 = 3c = cF_4.
\]
Since both $(U_n)$ and $(cF_{n+2})$ satisfy the same Fibonacci recurrence and agree at $n=1,2$, induction gives $U_n=cF_{n+2}$. Subtracting $2c$ yields
\[
  \Delta_n = c(F_{n+2}-2).
\]
Summing from $i=1$ to~$n$ gives
\[
  \tau_n
  =
  c\sum_{i=1}^n (F_{i+2}-2)
  =
  c(F_{n+4}-2n-3),
\]
using the standard Fibonacci summation identity.
\end{proof}

\begin{cor}[Asymptotic inflation factor]
\label{cor:phi}
Under the bootstrap and uniform-payload hypotheses of \autoref{cor:fibonacci}, the structural inflation factor
\[
  \Phi_n := \frac{\Delta_n}{\Delta_{n-1}}
\]
converges to the Golden Ratio
\[
  \varphi = \frac{1+\sqrt{5}}{2}.
\]
\end{cor}

\begin{proof}
Under the bootstrap assumptions of \autoref{cor:fibonacci},
\[
  \Phi_n
  =
  \frac{\Delta_n}{\Delta_{n-1}}
  =
  \frac{F_{n+2}-2}{F_{n+1}-2}.
\]
The additive constants are asymptotically negligible relative to the Fibonacci growth, so the standard ratio convergence gives $\Phi_n \to \varphi$.
\end{proof}

\section{Cubical Agda mechanization}
\label{sec:mechanization}

\subsection{Mechanizing the theorem package}

For the purposes of the present paper, the relevant Agda code now splits into theorem-facing metatheory modules and the barrier-and-recurrence modules that they explain.
\begin{itemize}[nosep]
\item The theorem-facing modules are:
  \begin{itemize}[nosep]
  \item \path{Metatheory/Extensional.agda},
  \item \path{Metatheory/KanSubsumption.agda},
  \item \path{Metatheory/AdjunctionBarrier.agda}.
  \end{itemize}
\item On the recurrence side, the key modules are:
  \begin{itemize}[nosep]
  \item \path{Core/Nat.agda},
  \item \path{Core/AffineRecurrence.agda},
  \item \path{ObligationGraph/Recurrence.agda},
  \item \path{Test/Fibonacci.agda},
  \item \path{Test/MetatheorySmoke.agda}.
  \end{itemize}
\item On the sealing side, the key modules are \path{Saturation/Decomposition.agda}, \path{Saturation/AbstractionBarrier.agda}, and \path{Saturation/AbstractionBarrier9.agda}.
\end{itemize}
The metatheory modules implement the semantic dictionary used throughout the paper. In the extensional/UIP case an obligation contributes no new primitive cost because its semantic space collapses to a contractible one. In the cubical case the relevant notion is weaker and computational: an obligation contributes no new primitive cost when a uniform \texttt{hcomp}-style derivation synthesizes a witness from already exported boundary data.

\path{Metatheory/Extensional.agda} defines binary semantic coherence obligations as path-equality types $p \equiv q$ and proves \texttt{UIP-forces-depth-1} together with the paper-facing corollary \texttt{history-truncates-to-one}. This is the exact Agda realization of the extensional statement that UIP collapses binary coherence and forces coherence depth~$1$.

\path{Metatheory/KanSubsumption.agda} packages the open-box inputs consumed by cubical composition; instantiated at square types, this is exactly the missing-face problem for an arity-$3$ obligation. In particular it proves:
\begin{itemize}[nosep]
\item \texttt{arity3-obligation-syntactically-derivable},
\item \texttt{history-beyond-two-algorithmically-subsumed},
\item \texttt{arity3-open-box-hfilled}.
\end{itemize}
This is the mechanized form of the telescopic subsumption claim: once the visible unary and binary faces are fixed, cubical composition uniformly computes a witness for the missing arity-$3$ face. The mechanization therefore establishes inhabitation by a generic derivation scheme, not singletonhood of the full geometric filler space.

\path{Metatheory/AdjunctionBarrier.agda} supplies the lower bound. It constructs a local two-point type, its swap autoisomorphism, the induced nontrivial path \texttt{swap-path}, and the contradiction theorem \texttt{depth1-insufficient}. Thus the cubical setting does not admit a global depth-one collapse of binary coherence spaces.

\subsection{Barrier discipline and affine recurrence}

Its organizing invariant is that resolved layers are exported through opaque interfaces. At the level of proof architecture, this is the computational form of \autoref{thm:trace}: once an obligation has been discharged, later constructions depend only on the exported schema and not on the hidden derivation tree that discharged it. In the barrier modules, this is checked locally for the transitions around steps~8 and~9 by showing that inherited obligations are discharged from the owning layer's opaque record interface alone, without reopening deeper internal derivations.

This realizes the abstraction barrier promised by \autoref{def:sealed}. It is the reason historical depth is measured in terms of sealed layers rather than in terms of raw derivation ancestry.

On the counting side, the Agda development makes the modeling discipline explicit. The active primitive interface is represented as the disjoint union of the two most recent schema sets, and the saturation step is stated separately rather than hidden in the syntax of imports. The concrete decompositions in \texttt{Saturation/Decomposition.agda} then record, step by step, how the next obligation set splits into contributions from the two most recent sealed layers.

\path{Core/AffineRecurrence.agda} strengthens the old normalized counting story to the payload-aware law
\[
  \Delta_{n+1} = \Delta_n + \Delta_{n-1} + 2c,
\]
defines the shifted sequence $U_n=\Delta_n+2c$, and proves that $U$ satisfies the homogeneous Fibonacci recurrence. The zero-payload specialization is linked back to the original \path{Core/Nat.agda} development, while \path{ObligationGraph/Recurrence.agda} remains as the stable payload-free corollary surface and \path{Test/Fibonacci.agda} continues to check the concrete Fibonacci identities.

A useful feature of the mechanization is that it separates three claims that are easy to conflate in prose. The \emph{abstraction-barrier claim} says that historical memory is mediated only by opaque exports. The \emph{metatheoretic depth claim} says that cubical filling uniformly derives witnesses for arity-$3$ obligations while binary obligations remain genuinely independent in general. The \emph{recurrence claim} says that, once the depth-two window and explicit payload accounting are in place, the cardinality of the next primitive interface is governed by the affine depth-two law above. The current Agda development now addresses all three points.

\begin{rem}[Scope of the mechanization]
\label{rem:mech-scope}
The present Agda artifact formalizes both the foundational causes and the combinatorial consequences of the coherence window. It proves that UIP yields depth~$1$, that cubical arity-$3$ obligations are algorithmically subsumed by uniform composition, that binary coherence does not collapse globally in the cubical setting, and that explicit payload accounting yields the affine recurrence whose constant shift is Fibonacci. The mathematical arguments of \autoref{sec:coherence} remain important as explanation and as links to the categorical and topological examples, but they no longer stand in for an absent mechanization of the paper's main depth claims.
\end{rem}

\section{Conclusion}
\label{sec:conclusion}

The paper isolates \emph{coherence depth} as a metatheoretic invariant of framework extension in type theory. In extensional settings with UIP, higher coherences collapse and the invariant is~$1$. In the fully coupled cubical univalent setting studied here, the invariant is~$2$. The difference between those two cases is not cosmetic: under uniform payload size it changes the recurrence law of integration from one-step affine growth to shifted Fibonacci scaling.

The mechanism behind the recurrence is structural. For globally acting foundational extensions, maximal interface density, enforced as a canonicity requirement, identifies sealing cost with the size of the active interface. The Integration Trace Principle decomposes each new public interface into the layer's own payload together with the discharged obligations that make it cohere with prior exports. In the cubical univalent case, telescopic subsumption shows that all older comparisons are algorithmically synthesized from the adjacent exported traces, so the primitive active interface is generated by the previous two layers; under constant payload size this yields an affine depth-two recurrence whose shifted form is Fibonacci.

From this viewpoint the appearance of the Golden Ratio is not an empirical regularity of one implementation and not a matter of aesthetic numerology. It is the characteristic root of the homogeneous recurrence recovered after separating the constant payload contribution from the depth-two coherence memory.

The current Cubical Agda artifact now matches that story. It proves the extensional collapse, the cubical upper bound, the lower bound against depth~$1$, and the affine recurrence that shifts to Fibonacci. The next natural step is therefore not to formalize the depth theorem package itself, but to connect this metatheory layer more tightly to the repository's oracle and candidate-generation components.

\bibliographystyle{alphaurl}
\bibliography{coherence_depth_refs}

\end{document}
